\documentclass[11pt]{article}

\usepackage[margin=1in]{geometry}
\usepackage{amsmath, amssymb, amsthm}
\usepackage{mathtools}
\usepackage{bm}
\usepackage{xcolor}
\usepackage{hyperref}
\usepackage[most]{tcolorbox}
\usepackage{booktabs}
\usepackage{enumitem}
\usepackage{listings}

\hypersetup{colorlinks=true, linkcolor=blue!50!black, urlcolor=blue!50!black,
            unicode=true, bookmarksdepth=2}

\definecolor{codebg}{HTML}{F3F4F6}
\definecolor{accent}{HTML}{C85020}
\definecolor{navy}{HTML}{12264E}
\definecolor{ok}{HTML}{1F7A4C}
\definecolor{warn}{HTML}{B45309}

\lstdefinestyle{stata}{
  basicstyle=\ttfamily\small,
  backgroundcolor=\color{codebg},
  frame=none, breaklines=true, columns=fullflexible,
  keepspaces=true, xleftmargin=4pt, framexleftmargin=4pt,
  commentstyle=\color{gray!90!black}\itshape,
  keywordstyle=\color{navy}\bfseries,
  morekeywords={reghdfe, hdfe, gen, generate, replace, drop, capture,
                bys, bysort, gegen, levelsof, forvalues, foreach, if,
                local, scalar, matrix, predict, sum, total, accum, invsym,
                det, import, use, save, tempfile, syntax, varlist},
  morecomment=[l]{*},
  morecomment=[l]{//},
}
\lstset{style=stata}

\newcommand{\code}[1]{\texttt{\small #1}}
\newcommand{\Adoref}[1]{\textcolor{accent}{\footnotesize\textbf{[ado L#1]}}}

\newcommand{\R}{\mathbb{R}}
\newcommand{\E}{\mathbb{E}}
\newcommand{\1}{\mathbb{1}}
\newcommand{\Smpl}{\mathcal{S}}
\newcommand{\thetad}{\widehat{\bm{\theta}}_d}
\newcommand{\Ngt}{N_{g,t}}
\newcommand{\Pproj}{P^{\!W}}
\newcommand{\Mproj}{M^{\!W}}

\title{\textbf{HDFE controls in \texttt{did\_multiplegt\_dyn}}\\[3pt]
       \large A fast path that bypasses the per-control loop \\
       \large for dummy-basis controls like \texttt{ac\_uq\_id\#monthyear}}
\author{}
\date{}

\begin{document}
\maketitle

% ============================================================================
\section{What you are trying to do, and why it is slow today}

You want to control for \code{ac\_uq\_id\#monthyear} on the
\code{data\_test.csv} panel. \code{ac\_uq\_id} is a coarser group attribute
(time-invariant at the group level), \code{monthyear} is the time index, and
\code{ac\_uq\_id\#monthyear} produces one dummy per (cluster, month) cell.
On the size of your data this means \textbf{12{,}000+ dummies} ---
generated, first-differenced, residualized, fed to \code{matrix accum},
and then to a per-baseline \code{invsym} of a $12000 \times 12000$ matrix.

\begin{tcolorbox}[colback=codebg, colframe=warn, boxrule=0.6pt, arc=1pt,
                  title=\textcolor{warn}{\textbf{Why naive \code{controls()} is intractable here}}]
The existing path treats every control symmetrically:
\begin{itemize}[leftmargin=*, itemsep=0pt]
\item Stage 1, L902--941: generate one $\Delta X_k$ and one
      $\widetilde{\Delta X}_k$ per control.
\item Stage 1, L1018: \code{matrix accum} on $K+1$ residualized variables
      \,---\, a matrix of size $K \times K$ when $K = 12000$.
\item Stage 1, L1034: \code{invsym(didmgt\_XX)} \,---\, $O(K^3)$ in $K=12000$
      is roughly $10^{12}$ flops per baseline cell.
\item Stage 2, L4679: $K$ \code{replace} commands per horizon $\ell$.
\item Stage 3, L4905--4954: a nested $K^2$ loop on $K=12000$ is $\sim 10^8$
      iterations per horizon per baseline.
\end{itemize}
\textbf{This is the wrong tool for dummy-basis controls.} The right tool is
fixed-effect absorption, and it gives the \emph{same estimator}.
\end{tcolorbox}

% ============================================================================
\section{The math fact that powers the fast path: FWL}

Take any partition of your dataset's columns into ``numeric controls'' $X$
and ``dummy controls'' $Z$ (where $Z$ are the indicators of the cells
$\mathrm{ac\_uq\_id}\times \mathrm{monthyear}$). The OLS of $\Delta Y$ on
$\Delta X$ and $\Delta Z$ gives the same numeric-controls slope $\thetad$ as
the OLS of $\Delta Y$ on $\Delta X$ after \emph{partialling out the column
space of $\Delta Z$ from both sides}:

\begin{equation}
\Bigl[\underbrace{\Delta X\;\;\Delta Z}_{=:\,\mathbf{X}_{\text{full}}}\Bigr]^{\!\top}\!
\mathbf{W}
\mathbf{X}_{\text{full}}\!
\begin{bmatrix}\thetad \\ \widehat{\phi}_d\end{bmatrix}
\,=\, \mathbf{X}_{\text{full}}^{\!\top}\!\mathbf{W}\Delta Y
\quad\iff\quad
\thetad \,=\, \bigl((\Mproj_{Z}\!\Delta X)^{\!\top}\mathbf{W}(\Mproj_{Z}\!\Delta X)\bigr)^{\!-1}
(\Mproj_{Z}\!\Delta X)^{\!\top}\mathbf{W}\Mproj_{Z}\!\Delta Y,
\label{eq:fwl}
\end{equation}
where $\Mproj_{Z} = \mathbf{I} - \Pproj_Z$ is the (weighted) annihilator of
$\Delta Z$. Frisch--Waugh--Lovell, exactly as in any textbook.

\medskip\noindent
\textbf{Reading.}\;The estimator we want \emph{never has to materialize}
the 12{,}000 columns of $\Delta Z$. We only need their \emph{projection
action}: ``demean by the column space of $\Delta Z$.'' For
$Z = \1\{a_g = a_0\}\cdot \1\{t = t_0\}$, that action is computable in
$O(N)$ per iteration of the alternating-projection algorithm used by
\code{hdfe}/\code{reghdfe} \,---\, no $12000\times 12000$ inversion.

\begin{tcolorbox}[colback=codebg, colframe=ok, boxrule=0.6pt, arc=1pt,
                  title=\textcolor{ok}{\textbf{What the fast path is, in one line}}]
``HDFE controls'' = the dummy columns of $\Delta Z$ are routed through
\code{hdfe}'s projection instead of through \code{matrix accum}. By
\eqref{eq:fwl} the $\thetad$ on the genuine numeric controls is unchanged.
\end{tcolorbox}

% ============================================================================
\section{The two flavors you would pass, and what each does}

The first-difference structure of \code{did\_multiplegt\_dyn} would seem to
rule out dummies that do not vary over time \,---\, their FD is zero. But the
help file already gives the workaround:
\begin{quote}\itshape
``To control for time-invariant covariates, one can for instance input the
product of those covariates and of the time variable $T$ into the option.''
\end{quote}
We apply that trick to both flavors of HDFE the user might want to absorb.

\subsection*{(A) Time-invariant dummies: \code{i.ac\_uq\_id} times \code{monthyear}}

For a group-level attribute $a_g$ that does not change over $t$, the bare
dummy has zero first difference:
\begin{equation}
\Delta \mathbf{1}\{a_g = a_0\} \,=\, \mathbf{1}\{a_g = a_0\} - \mathbf{1}\{a_g = a_0\} \,=\, 0.
\label{eq:bareinvariant}
\end{equation}
\textbf{The fix is to multiply by time.}\;Define
$\widetilde{Z}_{a_0}(g,t) := \mathbf{1}\{a_g = a_0\}\cdot t$. Then
\begin{equation}
\Delta \widetilde{Z}_{a_0}(g,t) \,=\, \mathbf{1}\{a_g = a_0\}\cdot t \,-\, \mathbf{1}\{a_g = a_0\}\cdot (t-1) \,=\, \mathbf{1}\{a_g = a_0\}.
\label{eq:Ztime}
\end{equation}
The first difference is the original dummy itself \,---\, generically non-zero.
So the construction is admissible as a control after all.

\medskip\noindent
\textbf{Interpretation.}\;Including
$\widetilde{Z}_{a_0} = \mathbf{1}\{a_g = a_0\}\cdot t$ for every cluster
level $a_0$ amounts to allowing each cluster to have its own
\emph{linear time trend} in $Y$. The coefficient on $\widetilde{Z}_{a_0}$
is the slope of cluster $a_0$'s trend relative to the omitted category.
This is the \texttt{trends\_lin}-style assumption applied at the cluster
level rather than the group level. With $A$ clusters this is $A$ extra
columns, which is the column space we want \code{hdfe} to absorb.

\subsection*{(B) Cluster-by-time dummies: \code{i.ac\_uq\_id\#i.monthyear} times \code{monthyear}}

The same trick applies to a cluster-by-time indicator. The bare cell-dummy
$\mathbf{1}\{a_g = a_0, t = t_0\}$ has a first difference that is non-zero
in two adjacent periods,
\begin{equation}
\Delta \mathbf{1}\{a_g = a_0, t = t_0\} \,=\, \mathbf{1}\{a_g = a_0\}\cdot\bigl(\mathbf{1}\{t = t_0\} - \mathbf{1}\{t = t_0 - 1\}\bigr),
\label{eq:cellfd}
\end{equation}
which already qualifies as a usable control. But to keep the construction
\emph{symmetric} with case (A) and to give it a linear-trend interpretation
inside each (cluster, time) cell, multiply by time:
\begin{equation}
\widetilde{Z}_{a_0, t_0}(g,t) := \mathbf{1}\{a_g = a_0, t = t_0\}\cdot t.
\end{equation}
The first difference is then
\begin{equation}
\Delta \widetilde{Z}_{a_0, t_0}(g,t) \,=\, \mathbf{1}\{a_g = a_0\}\cdot\bigl(t\cdot \mathbf{1}\{t = t_0\} - (t-1)\cdot \mathbf{1}\{t = t_0 - 1\}\bigr),
\label{eq:cellfd_t}
\end{equation}
which is non-zero in exactly the same two periods as \eqref{eq:cellfd} but
weighted by $t_0$ and $t_0 - 1$ respectively. The two specifications span
the same column space when both $\widetilde{Z}_{a_0, t_0}$ for every
$(a_0, t_0)$ \emph{and} the cluster-only $\widetilde{Z}_{a_0}$ are included
\,---\, FD-of-time-times-dummy together with FD-of-dummy together saturate
``cluster-by-time'' variation in linear trends.

\medskip\noindent
\textbf{Stata syntax for this object.}\;Either of the following two
equivalent specs:
\begin{lstlisting}
* Cluster-specific linear trend + cluster x time interaction trend:
hdfe_controls( i.ac_uq_id##c.monthyear )
* Or, written out:
hdfe_controls( i.ac_uq_id  i.ac_uq_id#c.monthyear )
\end{lstlisting}
The \code{c.monthyear} part is what implements the ``times year'' trick:
\texttt{reghdfe} interprets \code{i.ac\_uq\_id\#c.monthyear} as the
interaction of each cluster dummy with the continuous time variable,
giving exactly $\widetilde{Z}_{a_0} = \mathbf{1}\{a_g = a_0\}\cdot
\texttt{monthyear}$ for each $a_0$. The leading \code{i.ac\_uq\_id} would
be the cluster fixed effect itself; in our FD setup that piece is
absorbed by group differencing automatically and can be omitted (or kept
in the absorb spec where \code{reghdfe} will detect the collinearity
without error).

\subsection*{Both flavors at once, in absorb syntax}

When you want both case (A) and case (B) jointly, the natural absorb
spec is:
\begin{lstlisting}
hdfe_controls( i.ac_uq_id##c.monthyear  i.ac_uq_id#i.monthyear )
\end{lstlisting}
where the first term ($\mathbf{1}\{a_g\}\!\cdot\!1$ and
$\mathbf{1}\{a_g\}\!\cdot\!t$) gives cluster fixed effects and
cluster-specific linear trends, and the second term gives full
cluster-by-time intercepts. The two together span the column space of any
``constant for $g$ within $(a_g, t)$''-style HDFE the user might want.

\medskip
\noindent
This is the column space that \code{hdfe}'s projection is built to demean
by, and the option \code{hdfe\_controls()} simply forwards the user's
\code{reghdfe}-style spec into the \code{absorb()} argument of one
\code{hdfe} call per baseline level.

% ============================================================================
\section{Strategy: a focused option \code{hdfe\_controls()}}

\begin{lstlisting}
did_multiplegt_dyn Y G T D, effects(L) [other options] ///
                  controls(X1 X2)                       ///
                  hdfe_controls(ac_uq_id#monthyear)
\end{lstlisting}

\noindent
\textbf{Semantics.}\;The varlist in \code{hdfe\_controls()} is interpreted
\emph{not} as numeric controls but as a recipe for the projection
$\Pproj_Z$. Each entry is a categorical variable or interaction
(\texttt{reghdfe}-style absorb syntax). For each baseline level $d$, we
partial out the joint column space of these dummies from both $\Delta Y$ and
each genuine numeric control $\Delta X_k$, then estimate $\thetad$ on the
residualized inputs.

\medskip
\noindent
\textbf{What this option does NOT do} (and why that is correct):
\begin{itemize}[leftmargin=*, itemsep=2pt]
\item \textbf{No long-difference subtraction at L4679 for the HDFE part.}
      L4679 subtracts $\thetad[k]\cdot \Delta X^{(\ell)}_{k,g,t}$ for each
      \emph{numeric} control $k$. The HDFE projection is already taken out
      of the comparison through the event-study estimator's structure: a
      switcher's $\Delta Y^{(\ell)}$ is compared against not-yet-switchers
      \emph{from the same control cells} via the demeaning at L4647--4664,
      whose denominator $\widehat{E}^{(d)}_{g,t}$ is now built with the
      HDFE-augmented projection. The HDFE long-difference is canceled by
      the comparison itself, not by a separate subtraction.
\item \textbf{No new variance correction $U^{\mathrm{var},X}$ for the HDFE.}
      $U^{\mathrm{var},X}$ in the companion paper captures the noise from
      estimating $\thetad$. We do not store ``one coefficient per HDFE
      dummy'' \,---\, we treat the projection as a nuisance subspace, the
      same way \code{trends\_nonparam} does. The asymptotics are the
      \texttt{trends\_nonparam} asymptotics with a richer projector, which
      Section 1.4 of dCdH (2024) covers verbatim for stratum$\times$time
      structures.
\end{itemize}

% ============================================================================
\section{Why this is the right answer mathematically}

Two facts make the fast path identical to ``include them as controls'':

\medskip
\noindent
\textbf{Fact 1 (FWL).}\;By \eqref{eq:fwl}, $\thetad$ on the genuine numeric
controls is the same whether you (i) include $\Delta Z$ as columns of the
design matrix or (ii) partial out $\Delta Z$ from both sides.

\medskip
\noindent
\textbf{Fact 2 (Variance).}\;In the dCdH (2024) framework, the asymptotic
variance with $K$ \emph{finite} numeric controls has the $U^{\mathrm{var},X}$
correction. With $K \to \infty$ (the HDFE limit), the asymptotic framework
treats the FE projection as a deterministic nuisance subspace and
\textbf{does not} add a $K$-growing variance correction \,---\, that is the
correct asymptotic posture for the absorbed dimensions. So when we route
12{,}000 dummies through \code{hdfe}, we are simultaneously (a) computing
the right point estimate by FWL and (b) using the right asymptotic variance
by the FE-projection asymptotics.

\medskip
\noindent
The \texttt{trends\_nonparam} option in the file is exactly this regime
applied to one stratum. We are generalizing to many.

% ============================================================================
\section{What changes in the .ado, by stage}

\subsection*{Stage 1 \,---\, the projection enlarges}

The projector that residualizes $\Delta Y$ and each $\Delta X_k$ becomes:
\begin{equation}
\Mproj_{\mathcal{F}_d}
\;,\quad
\mathcal{F}_d \;=\; \mathrm{span}\bigl\{\,\1\{t=t_0\},\ \1\{s_g = s_0\}\cdot \1\{t=t_0\},\ Z^{(p)}_{g,t}\,\bigr\},
\label{eq:Fd}
\end{equation}
where $Z^{(p)}$ ranges over the absorb terms in \code{hdfe\_controls()}.
Within each baseline cell $d$, one \code{hdfe} call demeans $\Delta Y$ and
all genuine $\Delta X_k$ by $\mathcal{F}_d$ in a single sweep.

\medskip\noindent
\textbf{Code touches.}\;The four substitutions are the same as in the
``cluster\_time\_fe'' plan, but for a different conceptual reason:
\begin{itemize}[leftmargin=*, itemsep=0pt]
\item L941: $\widetilde{\Delta X}_k = \sqrt{N_{g,t}}\cdot [\Mproj_{\mathcal{F}_d}\Delta X_k]$.
\item L955: $\widetilde{\Delta Y} = \sqrt{N_{g,t}}\cdot [\Mproj_{\mathcal{F}_d}\Delta Y]$.
\item L958: $\mathrm{prod\_X}_{k,g,t} = N_{g,t}\cdot [\Mproj_{\mathcal{F}_d}\Delta X_k]$.
\item L984--996: $\widehat{E}^{(d)}_{g,t} = [\Pproj_{\mathcal{F}_d}\Delta Y]$
      via \code{reghdfe diff\_y\_XX `controlsXX', absorb(`hdfe\_set') residuals(...)}.
\end{itemize}

\subsection*{Stages 2 and 3 \,---\, unchanged}

\begin{tcolorbox}[colback=codebg, colframe=ok, boxrule=0.6pt, arc=1pt,
                  title=\textcolor{ok}{\textbf{Zero downstream code changes}}]
Stages 2 and 3 read \code{coefs\_sq\_`d'\_XX} (numeric-controls slope only),
\code{inv\_Denom\_`d'\_XX} (built off the augmented projection),
\code{prod\_X*\_Ngt\_XX} (per numeric control only), and
\code{E\_y\_hat\_gt\_int\_`d'\_XX} (now includes the HDFE part). All four are
correctly rebuilt by the Stage-1 patch. The long-difference residualization
at L4679 and the variance correction at L4905--4954 then run verbatim and
produce the right event-study estimates and SEs.
\end{tcolorbox}

% ============================================================================
\section{Performance estimate}

\begin{tabular}{@{}lccc@{}}
\toprule
Approach & Storage & Time per baseline $d$ & Numerical risk \\
\midrule
\code{controls(12000 dummies)}  &
  $12000\,N$ cells\,$+$\,$144{\cdot}10^6$ matrix &
  $O(K^3) = 10^{12}$ flops &
  singular \code{didmgt\_XX} \\
\code{hdfe\_controls(ac\_uq\_id\#monthyear)} &
  $N$ cells per demeaned variable &
  $O(N \cdot \text{iters})$, iters $\le 50$ &
  none (alternating projections) \\
\bottomrule
\end{tabular}

\medskip\noindent
On a panel with $N = 200{,}000$, $K = 12{,}000$, and 4--5 baseline levels,
the naive path is unrunnable in reasonable time on standard hardware. The
HDFE path is seconds: \code{hdfe} typically converges in fewer than 30
iterations, each linear in $N$.

% ============================================================================
\section{Stata patch}

\subsection*{(A) Syntax + validation \Adoref{53\,/\,65--73}}

\begin{lstlisting}
* L53 -- add the option (note the string type to accept reghdfe absorb syntax):
syntax varlist(min=4 max=4 numeric) [if] [in] [, /* ... */ hdfe_controls(string)]

* After the gtools check at L65-73:
if "`hdfe_controls'" != "" {
    qui cap which reghdfe
    if _rc { di as error "hdfe_controls() requires reghdfe. Run: ssc install reghdfe" ; exit 198 }
    qui cap which hdfe
    if _rc { di as error "hdfe_controls() requires hdfe. Run: ssc install hdfe" ; exit 198 }

    * Warn if the user passes a purely time-invariant variable (no-op).
    foreach a of local hdfe_controls {
        capture confirm variable `a'
        if !_rc {
            tempvar sd_a
            bys `2': egen `sd_a' = sd(`a')
            sum `sd_a'
            if r(max) == 0 {
                di as text "Note: `a' is time-invariant within group; it is absorbed by " ///
                           "differencing and will not contribute. Ignoring."
                local hdfe_controls : subinstr local hdfe_controls "`a'" "", word
            }
            drop `sd_a'
        }
    }
}
\end{lstlisting}

\subsection*{(B) Stage 1 \,---\, demean + theta + Den\_d (replaces L919--1098)}

\begin{lstlisting}
* Build the absorb spec: time FE, optional trends_nonparam x time, plus user's hdfe_controls.
local fe_set "time_XX"
if "`trends_nonparam'" != "" {
    capture drop trends_nonparam_temp_XX
    gegen trends_nonparam_temp_XX = group(`trends_nonparam')
    local fe_set "`fe_set' i.trends_nonparam_temp_XX#i.time_XX"
}
if "`hdfe_controls'" != "" {
    local fe_set "`fe_set' `hdfe_controls'"
}

* Allocate the residual containers used downstream.
local mycontrols_XX ""
local count_controls = 0
foreach var of varlist `controls' {
    local ++count_controls
    cap drop resid_X`count_controls'_time_FE_XX
    cap drop prod_X`count_controls'_Ngt_XX
    gen resid_X`count_controls'_time_FE_XX = 0
    gen prod_X`count_controls'_Ngt_XX = 0
    local mycontrols_XX "`mycontrols_XX' resid_X`count_controls'_time_FE_XX"
}
cap drop diff_y_wXX
gen diff_y_wXX = 0

levelsof d_sq_int_XX, local(levels_d_sq_XX)

foreach l of local levels_d_sq_XX {

    * Control sample for baseline d = `l'.
    tempvar smp
    gen byte `smp' = (ever_change_d_XX == 0       ///
                    & diff_y_XX != .              ///
                    & fd_X_all_non_missing_XX==1  ///
                    & d_sq_int_XX == `l'          ///
                    & time_XX < F_g_XX)

    * One hdfe call: demean DY and all numeric DX_k by F_d simultaneously.
    local dvars "diff_y_XX"
    forvalues k = 1/`count_controls' { local dvars "`dvars' diff_X`k'_XX" }

    qui hdfe `dvars' [aw=N_gt_XX] if `smp', ///
        absorb(`fe_set') generate(_resXX_)

    * Pack residuals into the names the unchanged tail expects.
    replace diff_y_wXX = sqrt(N_gt_XX) * _resXX_diff_y_XX ///
        if d_sq_int_XX == `l' & `smp' == 1
    forvalues k = 1/`count_controls' {
        replace resid_X`k'_time_FE_XX = sqrt(N_gt_XX) * _resXX_diff_X`k'_XX ///
            if d_sq_int_XX == `l' & `smp' == 1
        replace prod_X`k'_Ngt_XX = sqrt(N_gt_XX) * resid_X`k'_time_FE_XX ///
            if d_sq_int_XX == `l'
    }
    drop _resXX_*

    * --- E_y_hat_gt_int_l_XX: the full fitted value used downstream at L4664 ---
    cap drop E_y_hat_gt_int_`l'_XX
    local numericsXX ""
    forvalues k = 1/`count_controls' { local numericsXX "`numericsXX' diff_X`k'_XX" }

    cap reghdfe diff_y_XX `numericsXX' [aw=N_gt_XX] ///
        if d_sq_int_XX == `l' & time_XX < F_g_XX,  ///
        absorb(`fe_set') noconstant residuals(_res_DY_l_XX)
    scalar rc_XX = _rc
    if rc_XX == 0 {
        gen E_y_hat_gt_int_`l'_XX = diff_y_XX - _res_DY_l_XX ///
            if d_sq_int_XX == `l' & time_XX < F_g_XX
        drop _res_DY_l_XX
        scalar df_a_`l'_XX = e(df_a)
    }
    else {
        drop if d_sq_int_XX == `l'
        noi di "Baseline `l' dropped because of insufficient observations."
    }

    drop `smp'

    * The tail of the existing block (L1018-1098) runs verbatim from here:
    * matrix accum overall_XX = diff_y_wXX `mycontrols_XX' if d_sq_int_XX==`l' & ...
    * matrix coefs_sq_`l'_XX = invsym(didmgt_XX)*didmgt_Xy
    * matrix inv_Denom_`l'_XX = invsym(didmgt_XX) * r(sum) * G_XX
    * etc.
}
\end{lstlisting}

\noindent
The key efficiency point: \code{count\_controls} above is the number of
\emph{genuine numeric controls only} (e.g.\ 2), not 12{,}000. The 12{,}000
dummies enter solely through the \code{absorb(`fe\_set')} argument, which
\code{hdfe} handles with $O(N)$-per-iteration alternating projections
instead of an $O(K^3)$ matrix inversion.

\subsection*{(C) Stages 2 and 3 \,---\, NO changes}

The blocks at L4582--4689 (long-difference residualization for effects),
L5145--5184 (placebos), and L4905--4954 (variance correction) compile and
run as written. They consume the Stage-1 outputs which now incorporate
the HDFE projection.

% ============================================================================
\section{End-to-end example}

\begin{lstlisting}
import delimited using "data_test.csv", clear

* Numeric controls: X1 X2.
* HDFE: ac_uq_id (time-invariant -> auto-dropped) and ac_uq_id#monthyear.
did_multiplegt_dyn Y unique_small_grid_id monthyear D, ///
    effects(4) placebo(2)                             ///
    controls(X1 X2)                                   ///
    hdfe_controls(ac_uq_id#monthyear)                 ///
    cluster(ac_uq_id)
\end{lstlisting}

\noindent
What is being estimated, per baseline $d$:
\begin{equation}
\thetad \;=\; \arg\min_{\bm{\theta},\,\alpha^{(d)}}
\sum_{(g,t)\in\Smpl_d} \Ngt\Bigl(\Delta Y_{g,t} - \Delta X_{g,t}^{\!\top}\bm{\theta} - \alpha^{(d)}_{a_g,\,t}\Bigr)^{2},
\label{eq:est}
\end{equation}
which is identical to what \code{controls(X1 X2 i.ac\_uq\_id\#i.monthyear)}
\emph{would} compute, except that the 12{,}000-dimensional $\alpha^{(d)}$ is
absorbed instead of materialized.

% ============================================================================
\section{Sanity checks}

\begin{enumerate}[leftmargin=*, itemsep=2pt]
\item \textbf{Single-attribute equivalence.}\;
      \code{hdfe\_controls(i.ac\_uq\_id\#i.monthyear)} should give the same
      coefficients and SEs as \code{trends\_nonparam(ac\_uq\_id)}, because
      both produce the same projection $\Pproj_{\mathcal{F}_d}$ on
      \code{data\_test.csv}.
\item \textbf{Time-invariant drop.}\;
      \code{hdfe\_controls(i.ac\_uq\_id)} should print a notice
      ``time-invariant within group, ignoring'' and run as if no HDFE were
      passed. The point estimates and SEs should match the no-controls-FE
      command.
\item \textbf{FWL self-check.}\;On a small sub-sample where 12{,}000
      dummies actually fit in memory, \code{controls(X1 X2 d1 d2 ...
      d12000)} and \code{hdfe\_controls(i.ac\_uq\_id\#i.monthyear)} should
      give identical $\thetad$ on (X1, X2) up to floating-point precision.
\end{enumerate}

% ============================================================================
\section{Summary}

\begin{tcolorbox}[colback=codebg, colframe=navy, boxrule=0.6pt, arc=1pt,
                  title=\textcolor{navy}{\textbf{The whole proposal in one paragraph}}]
The user wants 12{,}000+ \code{ac\_uq\_id\#monthyear} dummies as controls.
Including them via \code{controls()} is computationally infeasible
($O(K^3)$ inverse), and time-invariant dummies inside that set are killed
by differencing anyway. The fix is a new option
\code{hdfe\_controls(spec)} accepting \code{reghdfe}-style absorb syntax.
It enlarges Stage 1's projector $\mathcal{F}_d$ with the user's HDFE
dimensions, uses one \code{hdfe} call per baseline level to demean
$\Delta Y$ and the genuine numeric controls, then runs the existing
\code{matrix accum}, $\thetad$, and $\mathrm{Den}_d^{-1}$ code unchanged.
Stages 2 and 3 (L4582--4689 and L4905--4954) need no edits: they consume
the rebuilt Stage-1 outputs. By Frisch--Waugh--Lovell the point estimates
match the limit of ``infinitely many dummy controls,'' and the variance
follows the \code{trends\_nonparam}-style nuisance-projection asymptotics
of dCdH (2024) Sec.\,1.4.
\end{tcolorbox}

\end{document}
